\textbf{Research abstract:} Traditional mechatronic system design has
predominantly adhered to a sequential, ad-hoc methodology in which the
physical structure and control are developed independently. In this
conventional paradigm, mechanical design decisions are typically
finalized prior to the integration of control architecture, frequently
overlooking the co-dependent interactions between structural dynamics
and control performance. Such segregated philosophy may leed in
suboptimal system behavior, as the interdependencies between mechanical
and control domains are not systematically addressed during its design
process. On the other hand, \emph{computational concurrent design}---or
\emph{co-design}---enables the parallel optimization of both mechanical
structure and control algorithms. This integrated methodology enhances
overall system performance and adaptability, particularly in
high-precision applications where the interaction between structure and
control is of paramount importance. By leveraging on the methods of
co-design, the proposed research framework addresses the limitations of
traditional design approaches: suboptimal system performance,
disturbance attenuation, suboptimal sensor and actuator placement,
lacking adaptability under operational conditions, and low lifetime.

This research proposes a dynamic topology optimization framework for
high-precision mechatronics and robotics. By extending density-based
methods to include temporal evolution, the framework enables concurrent
optimization of structure and control. Multi-indexed density variables
represent material distribution, actuator placement, and time-dependent
control inputs, facilitating efficient co-design. The approach aims to
advance the design and performance of next-generation mechatronic
systems.
